% !TeX TXS-program:bibliography = txs:///biber
\documentclass[11pt,twocolumn]{article}

\usepackage[T1]{fontenc}
\usepackage[utf8]{inputenc}
\usepackage{amsmath,amssymb}
\usepackage{hyperref}
\usepackage{standalone}
\usepackage{pgfplots}
\usepackage{pgfplotstable}
\usepackage{graphicx}
\usepackage{subcaption}
\usepackage{smartref}
\usepackage{cleveref}
\usepackage{tikz}
\usepackage[backend=biber,style=ieee]{biblatex}
\addbibresource{library.bib}


\usetikzlibrary{positioning,arrows.meta}

% custom commands
\newcommand{\subg}{\preccurlyeq}
\newcommand{\subgind}{\preccurlyeq_{\mathrm{ind}}}
\newcommand{\Forb}{\operatorname{Forb}}
\newcommand{\occ}[3]{\operatorname{occ}_{#2}(#1,#3)}
\newcommand{\toggle}[2]{#1\mathbin{\oplus}#2}
\newcommand{\edcost}[1]{\lvert #1\rvert}
\newcommand{\ParetoMin}{\operatorname{ParetoMin}}


\title{FGCGen: Generating Synthetic Graph Datasets for Structural Rule Discovery}
\author{Lukas Böhm}
\date{\today}

\begin{document}
\maketitle

\begin{abstract}
Understanding the reasoning behind decisions of Graph Neural Networks (GNNs) is crucial for trustworthy machine learning.
The novel field of Structural Rule Learning aims to output global and interpretable logical rules where current explainers often just yield a mask highlighting salient subgraphs for every single instance, often omitting crucial global information or model-wide patterns.
Existing graph benchmark datasets either lack a known ground-truth rule or are designed in a contrived way that makes them unsuitable for rule discovery.
We identify multiple such failure cases in all existing benchmark datasets and present the \emph{Forbidden Graph Characterization Generator} (FGCGen) as a solution.
Given forbidden subgraphs, optional imposed fragments, graph order, and a density band or degree cap, FGCGen generates connected valid graphs.
Negatives have to remain close to a tightly defined decision boundary to ensure identifiability of the dataset rule.
Finally we show empirically that most current models are unable to detect the defining patterns and existing explainers cannot reliably highlight the fundamental rule in nontrivial properties.
\end{abstract}

\section{Introduction} \label{sec:intro}
% intro
Many scientific questions about graphs concern a relationship between structure and an outcome.
In chemistry, certain mandatory functional groups or forbidden toxicophores determine the behaviour of a molecule or reaction. In other scenarios like finance and social networks, graph structures can reveal fraud or spam networks within a graph \cite{TODO}.
While predictive models like GNNs have shown great success on these classification tasks, they do not naturally reveal the learned relationship.
% Explainers
(Graph) Explainers are attempts on extracting the underlying mechanism that lead to a decision, usually by highlighting a salient subgraph per instance \cite{yingGNNExplainerGeneratingExplanations2019, luoParameterizedExplainerGraph2020, yuanExplainabilityGraphNeural2021} or across the dataset \cite{yuanXGNNModelLevelExplanations2020, yuanExplainabilityGraphNeural2023}. More sophisticated approaches also output conjunctions of multiple patterns or simple graph rules \cite{azzolinGlobalExplainabilityGNNs2023, ragnoLogicbasedSelfExplainableGraph2025, armgaanGRAPHTRAILTranslatingGNN, mullerGraphChefLearningRecipe}.
% TODO: further explain post-hoc and ante-hoc, as well as local vs global explainers? (similar to X-BAIC paper)
% rule recovery / structural graph learning (explanations don't match rules)
Recovering a general rule that classifies all graphs in the dataset is the main goal of structural graph learning.
The main mismatch between current explainers and rule discovery is that many evaluation methods rely on thresholding the subgraph explanations which defeats many features of a structurally defined property where no single element is inherently more important. \cite{proszewskaBXAICDatasetBenchmarking2025}

% desired rule format and desired dataset
% TODO: mention the boundary nearness requirement


Existing real-world datasets do not have a known structural basis, leading to conflicting explanations given by different explainers.
Synthetic benchmarks are attractive because their labels and expected explanations can be known by construction.
Simple datasets that are cheap to generate, such as BA-Shapes, Spurious-Motifs, etc. are commonly used in graph explainability research.
They are unsuitable for structural graph learning, as they frequently expose shortcuts that any machine learning model latches onto before focusing on the structural task (\ref{fig:shortcuts}) as shown in \cref{sec:failures}.

On the other end of the spectrum, extremal graph search and conjecturing build exact datasets for many properties to find potential invariants and growth laws.
Systems like \cite{devillezPHOEGHelpsObtain2019, davilaAutomatedConjecturingTxgraffiti2025, hansenAutoGraphiXAutomatedSystem2000} have already discovered many proven conjectures.
As the entire graph space needs to be enumerated including the expensive computations of properties like tree-width, those datasets are usually heavily constrained and not.

\begin{figure}[ht]
	\centering
	\includestandalone[width=\linewidth]{figures/positioning}
	\caption{Overview of graph learning approaches. As complexity and expressivety increase, the realizable graph order decreases}
	\label{fig:overview}
\end{figure}

The impreciseness of simple approaches and the limited size, lack of control and cost of the conjecturing systems leaves a gap for datasets that are large and precise enough for structural graph learning, while still being reasonably cheap to generate. The gap is visualized in \cref{fig:overview}.
Creating such data is hard because structural graph classes occupy highly constrained regions of graph space.
Blind (rejection) sampling can make valid graphs vanishingly rare, while independently generated class pools can be separated by density, degree, construction artifacts, or a wide edit margin.
To this end, we present FGCGen, a generator that can takes a basis that is defined through a basis of mandatory patterns from which either all have to appear together or only selected candidates from the set, and a set of forbidden patterns.
For each pattern both the induced and ordinary subgraph relation is allowed.
In the induced case, the forbidden basis is equivalent to specifying hereditary property through a forbidden graph characterization (FGC), hence the name FGCGen.
Additionally, we can guide the generation further by aiming towards a target density or degree cap.
Varying these values allows for trivial generation of useful OOD datasets.
In principle, the generator works by growing an initial graph vertex-by-vertex.
Since addition cannot destroy a pattern, we can embed the imposed patterns into the initial graph.
The growth process then only needs to avoid the forbidden patterns during growth, an operation that is monotone/upward-closed under both subgraph relations.
Above a certain graph order and density, valid graphs become exponentially rare.
To navigate the vast and hostile graph space, 
% TODO:  , beam search, MCMCM annearling
% Boundary crossing
% Released data


\paragraph{Contributions}
This paper makes four contributions.
First the shortcomings of existing datasets are defined and the requirements and interfaces for a reliable generator are defined.
Second, we develop FGCGen. A precise generator for properties defined through ordinary and induced subgraph patterns. This is released as a software package, together with a dataset suite of selected properties.
Third, we generate suitable boundary-close negatives.
Lastly, the datasets are evaluated on state-of-the-art GNNs and Explainers.
% TODO: mention also that we found witness floors, mandatory hubs (cographs) and overrigid classes?



\section{Related Work}

\subsection{Benchmarks for Graph Learning and Explanation}

Tree-Cycles, Tree-Grid, BA-Shapes, and BA-2Motifs are the most common and frequently used synthetic benchmark datasets in graph explainability research \cite{TODO the mentioned datasets}. They define a (binary) graph classification task based on the presence of motifs appended to a random background graph. 
% Background graph is not controlled so pattern might appear accidentally.
ShapeGGen substantially improves tis design with topology, homophily and a node classification target \cite{TODO}.


Another category of synthetic datasets probe the expressive power of GNNs: EXP presents pairs of 2-WL distinguishable graphs \cite{abboudSurprisingPowerGraph2021}, BREC and families of strongly regular graphs span even more WL-levels  \cite{wangEmpiricalStudyRealized2024, balcilarBreakingLimitsMessage2021}.
These ask whether an architecture can distinguish specific non-isomorphic pairs, while we are interested in detecting discriminative patterns.
The adjacent question of counting relevant patterns further separates GNN architectures based on their expressiveness \cite{chenCanGraphNeural2020}.

GraphWorld samples broad populations of synthetic tasks to study how model rankings change across generator parameters. Performance changed significantly once the synthetic graphs left the region of popular academic benchmarks \cite{TODO}.
Most benchmarks have no or only trivial edge or node labels that are at best added separately through tools like ALBTER \cite{TODO} or GenCAT \cite{maekawaGenCATGeneratingAttributed2023}.
SkyMap jointly generates topology and labelling, thereby eliminating common class imbalance and feature distribution issues \cite{wassingtonSkyMapGenerativeGraph2024}.

% TODO: some words on Cora, OGB, MUTAG, and other commonly used real datasets? ... not really, they don't have ground truth.


\subsection{Global Explanations}

% TODO: rewrite to cite azzolin 2025 and 2026 here!
As opposed to instance-wise local explanations, a global explainers tries to explain the entire dataset (or model) directly, which is necessary for rule recovery.
GraphFramEx argues that explanation evaluation must include necessity, sufficiency, stability, runtime, and simple baselines rather than only mask overlap \cite{TODO}.
Formal work further shows that minimal subgraph explanations correspond to logical prime implicants only under restricted conditions \cite{weinbauerPrimeImplicantExplanations2025}.
Thus high fidelity to a predictor is not equivalent to recovering the data rule. 
Model-level explainers usually try to disentangle a concept bottleneck layer to retrieve reusable recipes made from a conjunction of patterns \cite{azzolinGlobalExplainabilityGNNs2023, armgaanGRAPHTRAILTranslatingGNN, gengLogicXGNNGroundedLogical2026, ragnoLogicbasedSelfExplainableGraph2025} while other approaches directly try to distill the rule into a decision tree \cite{mullerGraphChefLearningRecipe, pluskaLogicalDistillationGraph2024}.

\subsection{Constrained Structural Generation}

In classical graph theory many hereditary properties admit constructive generators: co-trees for cographs \cite{TODO}, partial elimination orderings for chordal graphs \cite{TODO}, partial cube expansion for median graphs \cite{ovchinnikovPartialCubesStructures2008} and more general approaches \cite{iordanskiConstructiveGraphTheory2020}.
While those are convenient and fast (usually polynomial time), they lack statistical control and don't provide a way to expose suitable negatives.

Rejection based random graph sampling is simple and reliable but for many hereditary classes the valid members become exponentially rare, especially among larger densities and orders \cite{friezeIntroductionRandomGraphs2015}.
Modeling the generation as a chain of edge toggles can work using MCMC-based approaches, but is often plagued by bad mixing performance due to the non-ergodic nature of most properties \cite{TODO}. % TODO: wang landau, etc.. 
% Graphons describe limits of dense graph sequences rather than provide an exact finite-graph membership oracle \cite{lovaszLargeNetworksGraph2012}
Exponential random graph models can encode selected statistics, but forbidden patterns can only be approached asymptotically \cite{TODO}.

Learned Generators, including diffusion models, require exactly the training data we aim to generate \cite{vignacDiGressDiscreteDenoising2023}.

% TODO: graph grammars including Hyperedge-replacement grammars 

The most promising branch for precise property-adhering graph generation are SAT based approaches.
The SAT Modulo Symmetries framework 
\cite{kirchwegerSATModuloSymmetries2024, janotaBreakingSymmetriesQuantified2025, kirchwegerSATBasedGenerationPlanar2023, kirchwegerCoCertificateLearningSAT2023}


\section{Failure Modes for Synthetic Graph Datasets}\label{sec:failures}
% TODO: introduce notation here and then make the failure modes more formal!!!

\subsection{Notation}

Use finite, simple, undirected graphs. Write $G=(V,E)$, $n=|V|$, and $G[S]$ for an induced subgraph. Use $\mathcal G$ for the graph universe, $D\subseteq\mathcal G$ for the declared population domain, and $P:D\to\{0,1\}$ for the complete task rule.
% ideally D = \mathcal G$
Throughout this paper in-property graphs satisfy $P(G)=1$ and negatives satisfy $P(G)=0$.
To constraint the set of valid negative to be boundary-close, we require negatives to not leave $\mathcal{G}$, just the class defined by $P$.

The generator supports two containment relations. We say $H \subg G$ is the ordinary monomorphism: an injective edge-preserving map where additional host edges among mapped vertices are allowed. $H \subgind G$ defines the induced subgraph isomorphism, a map that also preserves nonedges.



% common dataset problems
\cite{faberWhenComparingGround2021}


\begin{figure*}[t]
	\centering
	\begin{subfigure}[t]{0.49\textwidth}
		\centering
		\resizebox{\linewidth}{!}{\includestandalone{figures/shortcuts_synthetic}}
		\caption{Synthetic benchmarks and controls.}\label{fig:syn}
	\end{subfigure}\hfill
	\begin{subfigure}[t]{0.49\textwidth}
		\centering
		\resizebox{\linewidth}{!}{\includestandalone{figures/shortcuts_real}}
		\caption{Real data and corrected reductions.}\label{fig:real}
	\end{subfigure}
	\caption{Structural-feature shortcut audit.}\label{fig:shortcuts}
\end{figure*}



\section{FGCGen}

\subsection{Rule Language}

Since we allow both induced and ordinary containment relations, we define a presence feature:
$$
A_{H,R}(G)=
\begin{cases}
	\mathbf 1[H\subg G],&R=\mathrm{sub},\\
	\mathbf 1[H\subgind G],&R=\mathrm{ind}.
\end{cases}
$$

The allowed language of the generator is a CNF over the patterns $\mathcal{F} =\{(F_i,R_i)\}_i$ and $\mathcal{H}=\{(H_j,T_j)\}_j$:
$$
P(G)=1 \iff \bigwedge_i\neg A_{F_i,R_i}(G) \land \bigwedge_j A_{H_j,T_j}(G), \qquad G\in D.
$$
This can be trivially unrolled to a fragment of existential-quantified first order logic on graphs.
% TODO: the all vs any thing in the imposed side ....

The forbidden patterns alone define $P_{\Forb}=\bigwedge_i\neg A_{F_i,R_i}(G)$, a hereditary backbone, closed under induced vertex deletion, even with mixed matching relations.  % TODO: is this correct?

% To avoid expensive, yet redundant computations during generation we exploit the heredity and enforce an anti-chain amongst the foribdden patterns.





\section{Boundary Crossing}

Suitable negatives need to remain close to the boundary.
% TODO: positives are always boundary close, while negatives can be like $\varepsilon n^2$ far away \cite{shapiraTestingEstimationGraph2024, martinEditDistanceFunction2013}

% TODO: also make sure that this is not only for induced subgraphs. (consider ordianry and even imposed fragments?)

Given an in property graph $G \in \mathcal{G}$, we aim to retrieve a graph $G' = G \oplus \delta \not\in \mathcal{G}$ created through an edit set $\delta \subseteq \binom{V(G)}{2}$
% TODO: proof that an edit set always exists

% The edit cost $| \delta |$ and the witness count are two conflicting quantities.

Unlike saturation theory this boundary crossing is 


\printbibliography

\appendix

\section{Proofs}\label{app:proofs}


\end{document}
